\documentclass[journal]{IEEEtran}

\usepackage{iftex}
\ifXeTeX
  \usepackage{fontspec}
  \usepackage{xeCJK}
  \xeCJKsetup{CJKspace=true}
  \defaultfontfeatures{Ligatures=TeX}
  \IfFontExistsTF{TeX Gyre Termes}{%
    \setmainfont{TeX Gyre Termes}%
  }{%
    \IfFontExistsTF{Times New Roman}{%
      \setmainfont{Times New Roman}%
    }{%
      \setmainfont{STIX Two Text}%
    }%
  }
  \IfFontExistsTF{NanumBarunGothicOTF}{%
    \setCJKmainfont[
      ItalicFont={NanumBarunGothicOTF},
      SmallCapsFont={NanumBarunGothicOTF}
    ]{NanumBarunGothicOTF}%
    \setCJKsansfont{NanumBarunGothicOTF}%
    \setCJKmonofont{NanumBarunGothicOTF}%
  }{%
    \IfFontExistsTF{Noto Sans CJK KR}{%
      \setCJKmainfont[
        ItalicFont={Noto Sans CJK KR},
        SmallCapsFont={Noto Sans CJK KR}
      ]{Noto Sans CJK KR}%
      \setCJKsansfont{Noto Sans CJK KR}%
      \setCJKmonofont{Noto Sans CJK KR}%
    }{%
      \setCJKmainfont[
        Scale=0.92,
        ItalicFont={Apple SD Gothic Neo},
        SmallCapsFont={Apple SD Gothic Neo}
      ]{Apple SD Gothic Neo}%
      \setCJKsansfont[Scale=0.92]{Apple SD Gothic Neo}%
      \setCJKmonofont[Scale=0.92]{Apple SD Gothic Neo}%
    }%
  }
  % NanumBarunGothic's CJK OpenType table does not expose U+00B7.
  % Classify it as Latin punctuation so the main Latin font renders it.
  \xeCJKDeclareCharClass{Default}{"00B7}
\else
  \PackageError{TRACE-RPG}{이 한글 원고는 XeLaTeX으로 빌드해야 합니다}{}
\fi

\usepackage{amsmath,amssymb}
\usepackage{booktabs}
\usepackage{graphicx}
\usepackage{tabularx}
\usepackage{url}
\usepackage[hidelinks]{hyperref}

\newcommand{\system}{TRACE-RPG}

% Shared table rhythm; paragraph columns wrap instead of reducing type size.
\setlength{\tabcolsep}{4pt}
\renewcommand{\arraystretch}{1.20}

\title{Traceability 보충자료}
\author{}

\begin{document}
\maketitle

% 본문 표 번호를 보존한다: Table XIV.
\setcounter{table}{13}

\section{제안된 traceability instrument (미측정)}
{\raggedright
제안된 평가 도구(instrument)의 구성 개념은 관찰 가능한 요구사항, 산출물, 행위, 결과 사이에서 복원할 수 있는 대응 관계다. 이 도구는 정확성, 내부 추론 과정, 개발자 역량, 플레이어 효익을 측정하지 않는다. 각 차원 $d$에 대해 특정 실험 조건에 치우치지 않는 관찰자 목록(inventory) $U_d$와 각 단위 $u$의 필수 연결 $L_d(u)$를 사전에 동결한다. $N_d$는 $U_d$에 속한 적용 가능 단위 수이고, $n_d$는 그중 모든 필수 연결이 검증된 단위 수다. 독립 평가자가 정확한 원천과 버전, 모든 원시 근거 관계, 위치 식별자(locator)를 확인한 단위에만 단위 점수 1을 부여한다. 근거가 누락되었거나 접근 불가, 구버전, 상충, 미검증 상태인 경우에는 사유를 기록하고 0으로 집계한다. 목록이 완전하고 $N_d>0$일 때만 $n_d/N_d$를 보고하며, 평가자 간 불일치율과 인위적 손상 주입 대조군(seeded-corruption control) 결과도 함께 제시하되 가중 복합 점수는 산출하지 않는다. 목록이 불완전하면 INDETERMINATE로 판정하고 관측된 건수와 누락 범위를 함께 보고한다. N/A는 비적용성이 독립된 근거로 확인된 경우에만 허용한다. 수리 항목은 실제 수리 시도가 0건이었음이 독립적으로 확인된 경우에만 충족으로 인정하며, 비어 있는 로그만으로는 근거가 되지 않는다. PASS, FAIL, NOT\_RUN, STAGE\_NOT\_REACHED, OUT\_OF\_SCOPE, MISSING 상태는 엄격히 구분하며, \texttt{stage\_not\_reached}는 결코 PASS로 간주하지 않는다.
\par}
\begin{table}[t]
\caption{출처 기반 traceability 평가 차원(미측정 제안)}
\label{tab:traceability-rubric}
\centering
\footnotesize
\setlength{\tabcolsep}{4pt}
\begin{tabularx}{\columnwidth}{@{}l>{\raggedright\arraybackslash}X>{\raggedright\arraybackslash}p{0.27\columnwidth}@{}}
\toprule
\textbf{ID} & \textbf{필수 correspondence} & \textbf{단위 분자/분모} \\
\midrule
TR01 & task, 정확한 artifact revision/digest, activity, actor의 identity/activity binding & complete slot / required slot \\
TR02 & request/input $\rightarrow$ rule version $\rightarrow$ evaluated check/result $\rightarrow$ decision & complete chain / inventoried decision \\
TR03 & decision $\rightarrow$ 관측 state/artifact effect와 before/after 근거 & complete correspondence / outcome-required decision \\
TR04 & triggering failure $\rightarrow$ corrective package $\rightarrow$ exact revalidation $\rightarrow$ result & complete repair chain / actual repair attempt \\
TR05 & forward requirement $\rightarrow$ artifact $\rightarrow$ verification 및 reverse in-scope change mapping & forward/reverse를 별도 보고 \\
TR06 & non-author가 선언된 audit/replay target을 독립 실행 & exercised target / declared target \\
\bottomrule
\end{tabularx}
\par\vspace{2pt}\noindent\parbox{0.96\columnwidth}{\footnotesize 미측정 제안이다. score를 수집하거나 검증하지 않았다. Codebook/일차 출처: 데이터 및 코드 가용성.}
\end{table}

\end{document}
